\documentclass[12pt]{article}
\usepackage[T1]{fontenc}
\usepackage[utf8]{inputenc}
\usepackage{lmodern}
\usepackage{textcomp}
\usepackage{enumitem}
\usepackage{geometry}
\usepackage{graphicx}
\usepackage{amsmath, amssymb}
\geometry{margin=1in}
\begin{document}
\begin{center}
History - K-12 Level Assessment 2 \
Total Marks: 40 \
Answer all questions.
\end{center}

\section*{Multiple Choice (10 marks)}
\begin{enumerate}[label=\arabic*.]
\item Which tool technology is associated with Neandertals?
\begin{enumerate}[label=(\alph*)]
    \item Aurignacian
    \item Acheulean
    \item Mousterian
    \item both b and c
\end{enumerate}
\item Clovis points are NOT found in:
\begin{enumerate}[label=(\alph*)]
    \item Alaska.
    \item Mexico.
    \item Canada.
    \item the continental United States.
\end{enumerate}
\item Evidence from northwest coast societies indicates that what is ultimately needed for the development of social complexity is:
\begin{enumerate}[label=(\alph*)]
    \item agriculture.
    \item irrigation technology.
    \item social stratification.
    \item a food surplus.
\end{enumerate}
\item The most obvious material symbols of ancient state societies are:
\begin{enumerate}[label=(\alph*)]
    \item stone tablets.
    \item pyramids.
    \item monumental works.
    \item irrigation canals.
\end{enumerate}
\item Analysis of the teeth of a hominid species discovered in Malapa Cave in South Africa, dating just after 2 million years ago, indicates a diet that consisted primarily of:
\begin{enumerate}[label=(\alph*)]
    \item fruits, leaves, and C 3 plants.
    \item meat derived from carnivorous mammals.
    \item meat, fish, and some plants.
    \item C 4 grasses that may have included wild wheat.
\end{enumerate}
\end{enumerate}

\section*{True / False (10 marks)}
\textit{Answer True or False}
\begin{enumerate}[label=\arabic*.]
\setcounter{enumi}{5}
\item True or False: The first pharaohs emerged in Egypt around 5100 years ago.
\item True or False: Civilizations are characterized by food surpluses, specialists, urban settlements, and a system of record keeping.
\item True or False: Cuneiform is the earliest form of writing in the world, which developed in Mesopotamia.
\item True or False: By 1900 B.P., the Mayan culture did not have a sophisticated bronze producing industry.
\item True or False: Maize, beans, and squash formed the essential subsistence triad for many indigenous New World civilizations.
\end{enumerate}

\section*{Long Form Response (20 marks)}
\textit{Answer each question with a clear and complete explanation.}
\begin{enumerate}[label=\arabic*.]
\setcounter{enumi}{10}
\item Discuss the agricultural practices of the Mississippian culture, including their reliance on maize and squash, and analyze how these practices influenced their societal development.
\item Discuss the evidence of trade between the Indus Valley cities and ancient Mesopotamian city-states, highlighting the significance of Harappan and Mesopotamian seals in this context.
\end{enumerate}

\vfill
\noindent\textit{End of Paper}
\end{document}